% from ega1-3-fr.tex
\oldpage[I]{114}
En effet (\hyperref[I.3.3.11-fr]{3.3.11}), $X_{(S')}$ s'identifie au produit
$X\times_Y Y_{(S')}$ grâce au diagramme commutatif
\phantomsection\label{I.3.4.8.proof-diagram-fr}
\[
  \xymatrix{
    X\ar[d]_f & X_{(S')}\ar[l]_p\ar[d]^{f_{(S')}}\\
    Y & Y_{(S')}\ar[l]^q
  }
\]

En vertu de (\hyperref[I.3.4.7-fr]{3.4.7}), la relation $q(y')=f(x)$ pour
$x\in M$, $y'\in Y_{(S')}$ équivaut à l'existence d'un $x'\in X_{(S')}$ tel
que $p(x')=x$ et $f_{(S')}(x')=y'$, d'où le corollaire.

Le lemme (\hyperref[I.3.4.6-fr]{3.4.6}) se précise de la façon suivante~:

\begin{proposition}[3.4.9]
\label{I.3.4.9-fr}
Soient $X$, $Y$ deux $S$-préschémas, $x$ un point de $X$, $y$ un point de $Y$,
au-dessus du même point $s\in S$. L'ensemble des points de $X\times_S Y$ ayant
pour projections $x$ et $y$ est en correspondance biunivoque canonique avec
l'ensemble des types d'extensions composées de $k(x)$ et $k(y)$ considérés
comme extensions de $k(s)$ (Bourbaki, \emph{Alg.}, chap. VIII, \S~8, prop. 2).
\end{proposition}

Soit $p$ (resp. $q$) la projection de $X\times_S Y$ dans $X$ (resp. $Y$) et
soit $E$ le sous-espace $p^{-1}(x)\cap q^{-1}(y)$ de l'espace sous-jacent à
$X\times_S Y$. Notons d'abord que les morphismes
$\operatorname{Spec}(k(x))\to S$ et $\operatorname{Spec}(k(y))\to S$ se
factorisent en
$\operatorname{Spec}(k(x))\to\operatorname{Spec}(k(s))\to S$ et
$\operatorname{Spec}(k(y))\to\operatorname{Spec}(k(s))\to S$~; comme
$\operatorname{Spec}(k(s))\to S$ est un monomorphisme
(\hyperref[I.2.4.7-fr]{2.4.7}), il résulte de
(\hyperref[I.3.2.4-fr]{3.2.4}) que l'on a
\phantomsection\label{I.3.4.9.tensor-product-fibre-fr}
\[
  P=\operatorname{Spec}(k(x))\times_S\operatorname{Spec}(k(y))
  =\operatorname{Spec}(k(x))\times_{\operatorname{Spec}(k(s))}
   \operatorname{Spec}(k(y))
  =\operatorname{Spec}\bigl(k(x)\otimes_{k(s)}k(y)\bigr).
\]

Nous allons définir deux applications $\alpha:P_0\to E$, $\beta:E\to P_0$
réciproques l'une de l'autre ($P_0$ désignant l'ensemble sous-jacent au
préschéma $P$). Si $i:\operatorname{Spec}(k(x))\to X$ et
$j:\operatorname{Spec}(k(y))\to Y$ sont les morphismes canoniques
(\hyperref[I.2.4.5-fr]{2.4.5}), on prendra pour $\alpha$ l'application des
espaces sous-jacents correspondant au morphisme $i\times_S j$. D'autre part,
tout $z\in E$ définit par hypothèse deux $k(s)$-monomorphismes
$k(x)\to k(z)$ et $k(y)\to k(z)$, donc un $k(s)$-monomorphisme produit
tensoriel $k(x)\otimes_{k(s)}k(y)\to k(z)$ et par suite un morphisme
$\operatorname{Spec}(k(z))\to P$~; $\beta(z)$ sera l'image de $z$ dans $P_0$
par ce morphisme. La vérification du fait que $\alpha\circ\beta$ et
$\beta\circ\alpha$ sont les applications identiques découle de
(\hyperref[I.2.4.5-fr]{2.4.5}) et de la définition du produit
(\hyperref[I.3.2.1-fr]{3.2.1}). On sait enfin que $P_0$ est en correspondance
biunivoque avec l'ensemble des types d'extensions composées de $k(x)$ et
$k(y)$ (Bourbaki, \emph{Alg.}, chap. VIII, \S~8, prop. 1).

\subsection{Surjections et injections}
\label{subsection:I.3.5-fr}

\begin{env}[3.5.1]
\label{I.3.5.1-fr}
De façon générale, considérons une propriété $\mathbf{P}$ de morphismes de
préschémas, et les deux propositions suivantes~:
\begin{enumerate}
  \item[(i)] Si $f:X\to X'$, $g:Y\to Y'$ sont deux $S$-morphismes possédant
  la propriété $\mathbf{P}$, $f\times_S g$ possède la propriété $\mathbf{P}$.
  \item[(ii)] Si $f:X\to Y$ est un $S$-morphisme possédant la propriété
  $\mathbf{P}$, tout $S'$-morphisme $f_{(S')}:X_{(S')}\to Y_{(S')}$, déduit
  de $f$ par extension du préschéma de base, possède la propriété
  $\mathbf{P}$.
\end{enumerate}

